\section{引言}
\label{sec:intro}

维护脆弱的用户界面（User Interface，UI）测试套件是现代 Web 工程中持续存在的维护成本。一次层叠样式表（Cascading Style Sheets，CSS）类名重命名、一次 React 组件的重新嵌套，或是一次行为不变的处理函数重构，都足以让数十个端到端（End-to-End，E2E）Playwright 测试断裂。这一问题在图形用户界面（Graphical User Interface，GUI）测试录制重放与自动生成的研究综述中被多次列为关键挑战~\cite{li2022androidReplay, wang2025mobileGui}。无论是工业界工具（如 Healenium~\cite{healenium}）还是日益增长的学术工作~\cite{joseph2026zerocost, mouraReproBreak2026}，给出的自然回应都是\emph{自愈}: 在定位器失效时，自动将其重写为指向新文档对象模型（Document Object Model，DOM）中等价元素的形式。近年来，大语言模型（LLM）被广泛引入软件测试与修复任务~\cite{xiang2025llmRepair, li2025llmFuzzing}，自动测试修复（含定位器修复）是其中最直接受益的方向之一。

这一前景看似乐观，近期文献也使该问题看起来接近解决。Joseph~\cite{joseph2026zerocost} 通过一个确定性可访问性阶梯在单一 demo 上报告了 100\% 通过率；当代基于 LLM 的定位器修复工作通常借助一致性守门规则报告修复成功率或解释一致性。然而这些指标共享一处结构性盲区：它们衡量的是\emph{打过补丁的测试是否通过}，而非\emph{补丁是否指向开发者真正期望的元素}。一个总能找到“某个”可解析选择器（任何属性近似的元素皆可）的修复器存在隐蔽的危险性——每一次“通过但错指”的修复都会掩盖一次本应被作为回归暴露的真实 UI 变更。

我们在 ReproBreak~\cite{mouraReproBreak2026} 数据集中的真实 Playwright 失效用例上直接测量这种失败模式。在 koenig 的 75 个可达失效用例上，结果令人不安:

\begin{itemize}
\item \textbf{F1 —— 隐性误修复是默认行为，软提示只提供有限保护。}当真实目标元素被强制从候选集中移除时，朴素本地修复器（\texttt{qwen2.5-coder:7b}~\cite{qwenCoderTech2024}）在 \textbf{78.7\%} 的用例中给出语法正确但元素错误的选择器（59/75）；追加“如无候选契合则放弃”守门规则后逐用例判定完全不变。为应对模型单一带来的有效性威胁，本文进一步在同一 75 个用例和同一候选移除协议上加入 \texttt{gpt-4o} 对照：其误修复率从 Vanilla 的 84.0\%（63/75）降至 Abstain 的 68.0\%（51/75），但仍有多数用例会被静默改写到错误元素。
\item \textbf{F2 —— 两个低开销确定性杠杆带来零回归的局部改进。}HealReact 的 L1 静态基底在 koenig 上可达 80.6\%（75/93）的真实失效目标；在可达用例上，L3 修复器的 top-1 选择器精确匹配率为 77.3\%（58/75）。testId 加权检索与强锚点后置过滤经二因子析因分离后合计 +3 个新正确用例且零回归，但在 $n=75$ 下 McNemar 未达显著，因此本文仅将其作为辅助性工程经验。
\item \textbf{F3 —— 跨应用泛化受锚点文化差异的制约。}在另一 React+Playwright 代码库（\texttt{payloadcms/payload}，319 个失效配对）上，由于其采用 BEM 风格的模板字面量类名、当前静态抽取器无法对其求值，原始可达率从 koenig 的 80.6\% 跌至 3.8\%。本文将其诊断为一处明确且可修复的盲区，而非对方法本身的否证。
\end{itemize}

\noindent F1 为无人值守自愈流水线引入行为重放验证层提供了直接的经验性动机：即便更强模型能从软提示中获益，提示文本仍不足以成为可审计的安全边界。F3 给出一条可证伪的跨应用泛化论断：任何仅在单一应用上完成基准测试的定位器修复工具，都很可能高估其泛化能力，因为锚点文化（\texttt{data-testid} vs \texttt{data-kg-*} vs BEM 模板字面量）在不同 React 代码库之间存在显著差异。

此外，我们使用 HealReact 的 L1/L3 静态基底作为诊断测量平台，而非声称已经完成完整自愈系统。该基底包含一个抽象语法树（Abstract Syntax Tree，AST）抽取器和一个带确定性后置校准的 LLM 派生意图标注器，静态可达 koenig 93 个失效目标中的 80.6\%；在其之上，一个小型本地 L3 修复器在可达用例上正确回答了 77.3\%（在完整数据集上为 62.4\% 的静态修复代理指标\footnote{本文所有“静态修复代理指标”均指选择器解析器层面的元素一致性，不是端到端 Playwright 通过率；二者可能因 visibility、timing 与运行时 DOM 差异而不一致。}）。两个低开销确定性杠杆 —— testId 加权检索与强锚点后置过滤 —— 经二因子析因分离后合计 3 个新正确用例、零回归（方向一致，McNemar 未达显著），表明确定性后处理可作为比提示工程更可审计的安全网。

\subsection{贡献}
（i）在真实 Playwright 数据集上对隐性误修复进行选择器层面的诊断性测量，并展示软提示放弃指令在本地 7B 模型上逐用例无效、在 \texttt{gpt-4o} 上仍留下高比例误修复（\S\ref{sec:findings}）；
（ii）使用 HealReact 的 L1/L3 静态基底作为可复现测量平台，而非声称已实现完整 L1--L4 自愈系统（\S\ref{sec:l1}）；
（iii）给出两个辅助性工程观察：低开销确定性杠杆可带来零回归的局部改进但统计未显著，跨应用可达率受锚点文化差异显著制约（\S\ref{sec:findings}）。

78.7\% 是选择器层面的对抗性上界（对抗性构造）；将其与 koenig 真正不可达的 18/93 用例朴素相乘，得到无人值守工作负载下约 15\% 的\emph{粗略上界估计}（而非测得的部署率，详见 \S\ref{sec:findings}）。即便作为上界估计，我们（\S\ref{sec:discussion}）仍认为这已经高到足以拒绝在持续集成（Continuous Integration，CI）中无人值守地自动提交 LLM 提出的选择器补丁。
